\section{Letter Combinations of a Phone Number}
\label{sec:rec-letter-combinations}

\problemstatement{Given a string of digits from $2$ to $9$, where each
digit maps to a fixed set of letters exactly as on a telephone keypad
(e.g.\ $2 \to \texttt{abc}$, $3 \to \texttt{def}$, \ldots, $7 \to
\texttt{pqrs}$, $9 \to \texttt{wxyz}$), return every possible letter
combination obtainable by choosing one letter for each digit, in order.}

\begin{patternbox}
\emph{``Choose one option per position, from a different choice set at
each position''} generalizes Section~\ref{sec:rec-power-set}'s binary
pick/not-pick into an $n$-ary branch at every level: instead of $2$
choices (include/exclude) at each position, there are however many
letters the current digit maps to, and the recursion simply loops over
all of them instead of hard-coding exactly two branches.
\end{patternbox}

\subsection*{Understanding the problem carefully}

Each digit position contributes one letter to the final combination,
chosen independently from that digit's fixed mapping (e.g.\ digit
$\texttt{'2'}$ contributes one of $\texttt{a},\texttt{b},\texttt{c}$).
Since these choices are independent across positions, exactly as the
pick/not-pick choices were independent across array elements in
Section~\ref{sec:rec-power-set}, the total number of combinations is the
\textbf{product} of each digit's letter count, and a recursion that
loops through all letters for the current digit, recursing once per
letter, visits every combination exactly once.

\begin{ideabox}[Loop Over the Current Digit's Letters Instead of a Fixed Binary Choice]
$\textsc{Solve}(\textit{index}, \textit{current})$: if
$\textit{index} = |\textit{digits}|$: record $\textit{current}$;
return. Otherwise: for each letter $c$ in the mapping for
$\textit{digits}[\textit{index}]$: append $c$ to $\textit{current}$;
recurse $\textsc{Solve}(\textit{index}+1, \textit{current})$; remove $c$
(undo).
\end{ideabox}

\subsection*{Why looping generalizes the binary pick/not-pick template directly}

The pick/not-pick recursion in Section~\ref{sec:rec-power-set} is the
special case of this loop where every position's choice set happens to
have exactly $2$ options (``include'' or ``exclude''). Here, the choice
set size varies by digit (e.g.\ $3$ letters for most digits, but $4$ for
$\texttt{7}$ and $\texttt{9}$), so the recursion is written as a general
loop over however many choices exist at the current position, rather
than two separately hard-coded branches -- the same choose/recurse/un-choose
rhythm applies uniformly, regardless of how many choices the loop
iterates over.

\subsection*{Algorithm}

\begin{enumerate}
  \item Define the digit-to-letters mapping (a fixed lookup table).
  \item $\textsc{Solve}(\textit{index}, \textit{current})$:
  \begin{enumerate}
    \item If $\textit{index} = |\textit{digits}|$: record
          $\textit{current}$; return.
    \item For each letter $c$ mapped from
          $\textit{digits}[\textit{index}]$: append $c$; recurse
          $\textsc{Solve}(\textit{index}+1, \textit{current})$; remove
          (undo).
  \end{enumerate}
  \item Call $\textsc{Solve}(0, \texttt{""})$ (return an empty list
        immediately if $\textit{digits}$ is empty, as a special case).
\end{enumerate}

\subsection*{Dry run}

\dryrun{Take $\textit{digits} = \texttt{"23"}$ ($2\to\texttt{abc}$,
$3\to\texttt{def}$).}

\begin{itemize}
  \item $\textsc{Solve}(0,\texttt{""})$: loop over $\texttt{a,b,c}$ for
        digit $\texttt{'2'}$.
  \begin{itemize}
    \item $c=\texttt{a}$: $\textsc{Solve}(1,\texttt{"a"})$: loop over
          $\texttt{d,e,f}$ for digit $\texttt{'3'}$: record
          \texttt{"ad"}, \texttt{"ae"}, \texttt{"af"}.
    \item $c=\texttt{b}$: $\textsc{Solve}(1,\texttt{"b"})$: record
          \texttt{"bd"}, \texttt{"be"}, \texttt{"bf"}.
    \item $c=\texttt{c}$: $\textsc{Solve}(1,\texttt{"c"})$: record
          \texttt{"cd"}, \texttt{"ce"}, \texttt{"cf"}.
  \end{itemize}
\end{itemize}

Result: \texttt{"ad","ae","af","bd","be","bf","cd","ce","cf"} -- all
$3\times3=9$ combinations.

\subsection*{C++ implementation}

\begin{lstlisting}
#include <bits/stdc++.h>
using namespace std;

void solve(int index, string& digits, string& current,
           vector<string>& mapping, vector<string>& result) {
    if (index == (int)digits.size()) {
        result.push_back(current);
        return;
    }
    string letters = mapping[digits[index] - '0'];
    for (char c : letters) {
        current.push_back(c);
        solve(index + 1, digits, current, mapping, result);
        current.pop_back();
    }
}

vector<string> letterCombinations(string digits) {
    if (digits.empty()) return {};

    vector<string> mapping = {
        "", "", "abc", "def", "ghi", "jkl",
        "mno", "pqrs", "tuv", "wxyz"
    };

    vector<string> result;
    string current;
    solve(0, digits, current, mapping, result);
    return result;
}

int main() {
    for (string& s : letterCombinations("23")) cout << s << " ";
    cout << endl;
    return 0;
}
\end{lstlisting}

\begin{complexitybox}
\textbf{Time Complexity.} If each digit maps to at most $4$ letters,
there are at most $4^L$ combinations for a digit string of length $L$,
each costing $O(L)$ to build, giving
\[
  O(L \cdot 4^L).
\]
\textbf{Space Complexity.} $O(L)$ for the recursion depth.
\end{complexitybox}

\begin{takeawaybox}
A loop-based recursive branch (``try every choice in this position's
choice set'') is the natural generalization of a hard-coded binary
pick/not-pick, and is the more common shape for backtracking problems in
practice, since most real choice sets have more than two options. Every
remaining backtracking problem in this chapter
(Sections~\ref{sec:rec-rat-in-maze}--\ref{sec:rec-word-search}) uses this
loop-based shape, with the choice set determined by the problem's
specific structure (grid moves, board placements, partition points).
\end{takeawaybox}
